\subsection{Учебные ступени полёта}

\textbf{Предисловие.} Итак, мы готовы приступить к программированию квадрокоптера. В этом разделе последовательно пройдём ключевые этапы: от безопасной предстартовой подготовки до взлёта, полёта по заданным точкам и посадки.

Любая автономная миссия дрона строится по стандартному жизненному циклу:

\begin{infobox}[Жизненный цикл миссии]
\begin{enumerate}
    \item \textbf{Инициализация:} Настройка таймеров, светодиодов, переменных.
    \item \textbf{Предстарт (Preflight):} Проверка систем. Дрон стоит на земле.
    \item \textbf{Взлёт (Takeoff):} Набор безопасной высоты.
    \item \textbf{Маршрут (Waypoints):} Основная часть программы. Полёт по точкам.
    \item \textbf{Возврат (RTL):} Возвращение в точку старта (опционально).
    \item \textbf{Посадка (Landing):} Снижение и выключение двигателей.
\end{enumerate}
\end{infobox}

Отдельной команды для «возврата» в API нет: этот этап — просто ещё один вызов \texttt{ap.goToLocalPoint(0, 0, z)}, возвращающий дрон в исходную точку сцены (см. §6) перед посадкой.

Каждый пример далее предваряется кратким пояснением, а строки кода снабжены комментариями, чтобы было понятно, что делает каждая команда.
